% Part: many-valued-logic
% Chapter: three-valued-logics
% Section: multiple-designation

\documentclass[../../../include/open-logic-section]{subfiles}

\begin{document}

\olfileid{mvl}{thr}{mul}

\olsection{يوازې $\True$ نۀ، نور قيمتونه هم ټاکل}

تر اوسه چې کوم منطقونه مو ليدلي، په ټولو کښې د ټاکل شوو
رښتياوالي قيمتونو سټ~$V^+ = \{\True\}$ و؛ يعنې يوه دعوه هله
رښتيا ګڼل کېده چې قيمت ئې~$\True$ و۔ خو کېدے شي هغه دعوه هم
رښتيا وګڼو چې يوازې~$\False$ نۀ وي۔ بيا د بېلګې په توګه په
ماتريس کښې~$V^+ = \{\True, \Undef\}$ ټاکل يو بل منطق جوړوي۔

\begin{defn}
\emph{د پارادوکس منطق}~$\LogLP$ په دې ماتريس تعريفېږي:
\begin{enumerate}
  \item معياري بياني ژبه~$\Lang L_0$ چې
  $\lnot$, $\land$, $\lor$, $\lif$
  نښلوونکي لري۔
  \item د رښتياوالي د قيمتونو سټ~$V = \{\True, \Undef, \False\}$۔
  \item $\True$ او~$\Undef$ دواړه ټاکل شوي دي؛ يعنې~$V^+ = \{\True, \Undef\}$۔
  \item صدق تابعې د کليني د قوي منطق له تابعو سره يو شان دي۔
\end{enumerate}
\end{defn}

\begin{defn}
د هالډېن \emph{د بې‌معنايۍ منطق}~$\LogHal$ په دې ماتريس تعريفېږي:
\begin{enumerate}
  \item معياري بياني ژبه~$\Lang L_0$ چې
  $\lnot$, $\land$, $\lor$, $\lif$
  او يو $1$ ځايه نښلوونکے~$+$ لري۔
  \item د رښتياوالي د قيمتونو سټ~$V = \{\True, \Undef, \False\}$۔
  \item $\True$ او~$\Undef$ دواړه ټاکل شوي دي؛ يعنې~$V^+ = \{\True, \Undef\}$۔
  \item صدق تابعې د کليني د کمزوري منطق له تابعو سره يو شان دي،
  خو د «بې‌معنا دے» عملګر هم ورسره شته:
  \begin{center}
    \begin{tabular}{c|c} 
    $\tf{+}$ & \\ 
    \hline  
    $\True$ & $\False$ \\ 
    $\Undef$ & $\True$ \\
    $\False$ & $\False$ 
    \end{tabular}
  \end{center}
\end{enumerate}
\end{defn}

دا منطقونه، سره له دې چې د کليني له منطقونو سره صدق جدولونه
شريکوي، د هغو پر خلاف \emph{تاوتولوژۍ لري}۔

\begin{prop}\ollabel{prop:LP-taut-CL} د~$\LogLP$ تاوتولوژۍ د کلاسيک
  بياني منطق له تاوتولوژيو سره يو شان دي۔
\end{prop}

\begin{proof}
  د~\olref[syn][sub]{prop:mvl-cl} له مخې، که~$\Entails[\LogLP] !A$
  وي، نو~$\Entails[\LogCL] !A$ هم وي۔ د برعکس لوري د ښودلو دپاره
  ثابتوو چې که داسې !!a{valuation}~$\pAssign v\colon \PVar \to \{\False, \True,
  \Undef\}$ وي چې~$\pValue v(!A)[\LogKs] = \False$، نو داسې
  !!a{valuation}~$\pAssign {v'}\colon \PVar \to \{\False, \True\}$
  هم شته چې~$\pValue {v'}(!A)[\LogCL] = \False$۔ دا د~$\LogLP$
  دپاره مطلوبه پايله ورکوي، ځکه~$\LogKs$ او~$\LogLP$ يو شان
  ځانګړې صدق تابعې لري، او~$\False$ د~$\LogLP$ يوازينے
  نۀ ټاکل شوے قيمت دے (همدا د~$\LogLP$ او~$\LogKs$ يوازينے توپير دے)۔
  نو که~$\Entails/[\LogLP] !A$ وي، د يوه~!!{valuation}~$\pAssign v$
  دپاره~$\pValue v(!A)[\LogLP] = \pValue
  v[\LogKs](!A) = \False$ وي۔ د ثابتېدونکې دعوې له مخې
  $\pValue{v'}[\LogCL](!A) = \False$، يعنې~$\Entails/[\LogCL] !A$۔
  
  د دې دعوې د ثبوت دپاره لومړے~$\pAssign {v'}$ داسې تعريفوو:
  \[
    \pAssign {v'}(p) = 
    \begin{cases}
    \True & \text{که } \pAssign {v}(p) \in \{\True, \Undef\}\\
    \False & \text{په بل حالت کښې}
  \end{cases}
  \]
  اوس پر~$!A$ په استقرا سره ښيو چې (الف) که~$\pValue v(!A)[\LogKs]
  = \False$ وي، نو~$\pValue {v'}(!A)[\LogCL] = \False$؛ او (ب) که
  $\pValue v(!A)[\LogKs] = \True$ وي، نو~$\pValue {v'}(!A)[\LogCL] =
  \True$۔
  \begin{enumerate}
    \item د استقرا بنسټ:~$!A \ident p$۔ د~\olref[syn][val]{defn:pValue}
    له مخې~$\pValue v(!A)[\LogKs] = \pAssign v(p)$ دے۔ که دا قيمت
    دروغ وي، تعريف شوې نوې ټاکنه هم دروغ ورکوي؛ او که رښتيا وي،
    نوې ټاکنه هم رښتيا ورکوي۔ له همدې امله (الف) او (ب) دواړه پوره دي۔
    د سرچينې يادونه: اصلي ثبوت دلته د دواړو ټاکنو بې‌قيده برابري
    ليکي؛ هغه د ناټاکلي مدخل دپاره ناسمه ده، خو دا دوه شرطي دعوې
    پرې نۀ ولاړېږي۔
    
    د استقرايي ګام دپاره لاندې حالتونه وګورئ:
    \item $!A \ident \lnot !B$. 
    \begin{enumerate}
      \item فرض کړئ~$\pValue v(\lnot!B)[\LogKs] = \False$۔ د
      $\tf{\lnot}[\LogKs]$ د تعريف له مخې~$\pValue v(!B)[\LogKs]  =
      \True$۔ د استقرايي فرض د (ب) حالت له مخې~$\pValue
      {v'}(!B)[\LogCL]  = \True$ کېږي؛ نو~$\pValue {v'}(\lnot !B)[\LogCL] = \False$۔
      \item فرض کړئ~$\pValue v(\lnot!B)[\LogKs] = \True$۔ د
      $\tf{\lnot}[\LogKs]$ د تعريف له مخې~$\pValue v(!B)[\LogKs]  =
      \False$۔ د استقرايي فرض د (الف) حالت له مخې~$\pValue
      {v'}(!B)[\LogCL]  = \False$ کېږي؛ نو~$\pValue {v'}(\lnot !B)[\LogCL] = \True$۔
    \end{enumerate}

    \item $!A \ident (!B \land !C)$.
    \begin{enumerate}
      \item فرض کړئ~$\pValue v(!B \land !C)[\LogKs] = \False$۔ د
      $\tf{\land}[\LogKs]$ د تعريف له مخې يا~$\pValue v(!B)[\LogKs]  =
      \False$ وي، يا~$\pValue v(!C)[\LogKs]  = \False$ وي۔ د
      استقرايي فرض د (الف) حالت له مخې يا~$\pValue {v'}(!B)[\LogCL]  = \False$
      وي، يا~$\pValue {v'}(!C)[\LogCL]  = \False$؛ نو~$\pValue {v'}(!B \land
      !C)[\LogCL] = \False$۔ د سرچينې يادونه: اصلي متن په دويم بديل
      کښې په تېروتنه بيا د لومړي فرعي فارمول قيمت ليکي؛ دلته د دويم
      فرعي فارمول قيمت ورکړل شو۔
      \item فرض کړئ~$\pValue v(!B \land !C)[\LogKs] = \True$۔ د
      $\tf{\land}[\LogKs]$ د تعريف له مخې~$\pValue v(!B)[\LogKs]  =
      \True$ او~$\pValue v(!C)[\LogKs]  = \True$۔ د
      استقرايي فرض د (ب) حالت له مخې~$\pValue {v'}(!B)[\LogCL]  = \True$
      او~$\pValue {v'}(!C)[\LogCL]  = \True$؛ نو~$\pValue {v'}(!B \land
      !C)[\LogCL] = \True$۔ د سرچينې يادونه: اصلي متن په دويم مدخل
      کښې په تېروتنه بيا د لومړي فرعي فارمول قيمت ليکي؛ دواړه مدخلونه
      بايد رښتيا وي۔
    \end{enumerate}
  \end{enumerate}

  نور دوه حالتونه ورته دي او تمرين ته پرېښودل شوي دي۔ په بله لار،
  پورته ثبوت د هغو ټولو~!!{formula}s دپاره پايله ثابتوي چې يوازې
  $\lnot$ او~$\land$ لري۔ بيا دا حقيقتونه کارولے شو چې په
  $\LogKs$ او~$\LogCL$ دواړو کښې د هر~$\pAssign v$ دپاره
  $\pValue v(!B \lor
  !C) = \pValue v(\lnot(\lnot!B \land \lnot !C))$ او~$\pValue v(!B \lif
  !C) = \pValue v(\lnot(!B \land \lnot !C))$ دي۔
\end{proof}

\begin{prob}
  د~\olref[mvl][thr][mul]{prop:LP-taut-CL} ثبوت بشپړ کړئ؛ يعنې
  په هغو حالتونو کښې (الف) او (ب) ثابت کړئ چې~$!A \ident (!B \lor !C)$
  او~$!A \ident (!B \lif !C)$ وي۔
\end{prob}

\begin{prob}
ثابته کړئ چې هره کلاسيکه تاوتولوژي په~$\LogHal$ کښې هم تاوتولوژي ده۔
\end{prob}

که څه هم د دغو منطقونو تاوتولوژۍ له کلاسيک منطق سره يو شان دي،
د استلزام اړيکې ئې بېلې دي۔ د بېلګې په توګه،~$\LogLP$
\emph{ناسازګاري زغمونکے} دے، ځکه چې~$\lnot p, p \Entails/ q$؛
نو د انفجار اصل~$\lnot !A, !A \Entails !B$ په عمومي ډول نۀ صدق کوي۔
(د~$!A$ او~$!B$ په ځينو حالتونو کښې صدق کوي، لکه چې~$!B$
تاوتولوژي وي۔)

\begin{prob}
  له لاندې اړيکو کومې (الف) په~$\LogLP$ او (ب) په~$\LogHal$
  کښې صدق کوي؟ د هرې يوې دپاره د صدق جدول ورکړئ۔
  \begin{enumerate}
    \item $p, p \lif q \Entails q$
    \item $\lnot q, p \lif q \Entails \lnot p$
    \item $p \lor q, \lnot p \Entails q$
    \item $\lnot p, p \Entails q$
    \item $p \Entails p \lor q$
    \item $p \lif q, q\lif r \Entails p \lif r$
  \end{enumerate}
\end{prob}

که په~$\LogLuk[3]$ کښې~$\Undef$ هم ټاکلے وګرځوئ، نو څه کېږي؟

\begin{defn}
  \emph{درې ارزښته آر-مينګل} منطق~$\LogRM[3]$ په دې ماتريس تعريفېږي:
  \begin{enumerate}
    \item معياري بياني ژبه~$\Lang L_0$ چې
    $\lfalse$, $\lnot$, $\land$, $\lor$, $\lif$ نښلوونکي لري۔
    \item د رښتياوالي د قيمتونو سټ~$V = \{\True, \Undef, \False\}$۔
    \item $\True$ او~$\Undef$ دواړه ټاکل شوي دي؛ يعنې~$V^+ = \{\True, \Undef\}$۔
    \item صدق تابعې د لوکاشېوېچ د منطق~$\LogLuk[3]$ له تابعو سره يو شان دي۔
  \end{enumerate}
\end{defn}

\begin{prob}
  له لاندې اړيکو کومې په~$\LogRM[3]$ کښې صدق کوي؟
  \begin{enumerate}
    \item $p, p \lif q \Entails q$
    \item $p \lor q, \lnot p \Entails q$
    \item $\lnot p, p \Entails q$
    \item $p \Entails p \lor q$
  \end{enumerate}
\end{prob}

ځينې وخت بېلابېل صدق جدولونه يوازې د ټاکل شوو قيمتونو په بدلولو
سره هماغه يو منطق (هماغه د استلزام اړيکه) جوړولے شي۔ د بېلګې په
توګه، د ګوډل په منطق کښې که د~$\{\True\}$ پر ځای~$V^+ = \{\True,
\Undef\}$ واخلو، همداسې کېږي۔

\begin{prop}\ollabel{prop:gl-udes}
  هغه ماتريس چې~$V = \{\False, \Undef,\True\}$، $V^+=\{\True,
  \Undef\}$ او د ګوډل د~$3$ ارزښته منطق صدق تابعې لري،
  کلاسيک منطق تعريفوي۔
\end{prop}

\begin{proof}
  تمرين۔
\end{proof}

\begin{prob}
  د~\olref[mvl][thr][mul]{prop:gl-udes} ثبوت ورکړئ۔ د دې دپاره
  هغه منطق~$\Log L$ وګورئ چې د ګوډل د منطق په شان تعريفېږي،
  خو~$V^+=\{\True,\Undef\}$ لري؛ وښيئ چې که
  $\Gamma \Entails/[\Log L] !B$ وي، نو~$\Gamma \Entails/[\LogCL] !B$
  هم وي۔ د~\olref[mvl][thr][mul]{prop:LP-taut-CL} د ثبوت فکرونه
  وکاروئ، خو د (الف) او (ب) ځانګړنو د ثبوت پر ځای وښيئ چې
  $\pValue v(!A)[\LogGod] = \False$ هله او يوازې هله وي چې
  $\pValue {v'}(!A)[\LogCL] = \False$ وي (او له دې امله~$\pValue
  v(!A)[\LogGod] \in \{\True,\Undef\}$ هله او يوازې هله وي چې
  $\pValue {v'}(!A)[\LogCL] =
  \True$ وي)۔ روښانه کړئ چې دا ولې قضيه ثابتوي۔
\end{prob}

\end{document}
